\documentclass[11pt,a4paper]{coverletter}

\senderblock
  {Radheem Bin Razi}
  {Distributed Systems \& Cloud-Native Engineer}
  {Ilmenau, Germany \textbar\ \href{mailto:sheikh.radheem@gmail.com}{sheikh.radheem@gmail.com} \textbar\ \href{https://radheem.github.io/my-cv}{portfolio}}

\begin{document}

%% ===== ENGLISH =====
\recipient{PwC Deutschland}{Hiring Team}{}

\opening{Dear Noemia Gryzia,}

Building infrastructure that quietly handles complex data flows while keeping security and compliance front and center is the part of backend engineering I value most. PwC’s Public \& Energy Transformation team tackles exactly those kinds of systemic challenges, and I want to bring that focus to the Backend-Developer role. My background centers on designing distributed systems that scale reliably, integrate securely across multiple vendors, and maintain strict data quality under real-world constraints.

At a recent platform project, I replaced fragile sidecar dependencies with a native NATS JetStream messaging layer and direct gRPC endpoints. This change reduced latency and simplified the deployment topology. I paired that architecture with a CUE-driven schema contract to validate incoming data from third-party vendors. The system also shipped a Ginkgo-based end-to-end data-quality gate that caught pipeline drift before it reached production. These patterns directly support the scalable, secure microservices and API design work your team delivers for public sector clients.

I also built a containerized event pipeline that fans out high-frequency metrics into Kafka. The pipeline routes those messages into time-series and document stores for real-time dashboards and long-term analytics. That work required a Ginkgo-based end-to-end data-quality gate, strict unit and integration tests, performance tuning under load, and careful credential management for every external connection. I apply Clean Code and SOLID principles consistently — at Bluefin Exchange I functioned as the system expert for code and design reviews, enforcing those standards across a distributed team before releases shipped. I am comfortable collaborating across business, technology, and operations teams to translate complex requirements into clean, maintainable code. PwC’s cross-functional delivery model aligns with how I prefer to work.

I am ready to contribute to your Public Sector \& Energy projects from day one. I can start full-time immediately. I am currently based in Germany and can relocate within the country within two to three weeks. My profile is Blue Card-ready, so a qualifying offer requires only a standard employment contract. Thank you for considering my application.

\closing{Sincerely,}{Radheem Bin Razi}

\clearpage
\selectlanguage{ngerman}

%% ===== DEUTSCH =====
\recipient{PwC Deutschland}{Personalabteilung}{}

\opening{Sehr geehrte/r Noemia Gryzia,}

Der Aufbau von Infrastruktur, die komplexe Datenflüsse im Hintergrund effizient verarbeitet, dabei aber Sicherheit und Compliance konsequent in den Vordergrund stellt, ist der Aspekt der Backend-Entwicklung, den ich am höchsten bewerte. Das Team Public \& Energy Transformation bei PwC widmet sich exakt diesen systemischen Herausforderungen, und ich möchte meine Expertise und diesen Fokus in die Position als Backend-Developer einbringen. Mein beruflicher Schwerpunkt liegt auf dem Entwurf verteilter Systeme, die zuverlässig skalieren, sicher über verschiedene Anbieter hinweg integriert werden können und unter realen Rahmenbedingungen höchste Datenqualität gewährleisten.

In einem aktuellen Plattformprojekt habe ich anfällige Sidecar-Abhängigkeiten durch eine native NATS JetStream-Messaging-Schicht sowie direkte gRPC-Endpunkte ersetzt. Diese Anpassung verringerte die Latenzzeiten und vereinfachte gleichzeitig die Deployment-Topologie. Ich ergänzte diese Architektur um einen CUE-basierten Schema-Contract, um eingehende Daten von Drittanbietern zu validieren. Zudem wurde im System eine Ginkgo-basierte End-to-End-Datenqualitätsprüfung implementiert, die Drift in der Datenpipeline erkennt, bevor diese die Produktionsumgebung erreichen. Diese Vorgehensweisen unterstützen unmittelbar die skalierbare und sichere Microservices- sowie API-Entwicklung, die Ihr Team für Kunden im öffentlichen Sektor umsetzt.

Darüber hinaus habe ich eine containerisierte Event-Pipeline entwickelt, die hochfrequente Metriken an Kafka weiterleitet. Die Pipeline leitet diese Nachrichten in Zeitreihen- und Dokumentendatenbanken weiter, um Echtzeit-Dashboards sowie langfristige Analysen zu ermöglichen. Für diese Aufgabe waren ein Ginkgo-basiertes End-to-End-Datenqualitäts-Gate, strenge Unit- und Integrationstests, Performance-Optimierungen unter Last sowie ein sorgfältiges Credential-Management für jede externe Verbindung erforderlich. Clean-Code- und SOLID-Prinzipien wende ich konsequent an — bei Bluefin Exchange fungierte ich als System-Experte für Code- und Design-Reviews und stellte sicher, dass diese Standards in einem verteilten Team vor jedem Release eingehalten wurden. Ich arbeite gerne und effektiv mit Fachabteilungen, Technologie- und Operations-Teams zusammen, um komplexe Anforderungen in sauberen, wartbaren Code zu übersetzen. Das crossfunktionale Delivery-Modell von PwC entspricht genau meiner bevorzugten Arbeitsweise.

Ich bin bereit, ab dem ersten Tag einen aktiven Beitrag zu Ihren Public Sector \& Energy-Projekten zu leisten. Ein sofortiger Vollzeitstart ist für mich möglich. Ich bin derzeit in Deutschland ansässig und kann innerhalb von zwei bis drei Wochen bundesweit umziehen. Mein Profil ist Blue Card-ready, sodass ein entsprechendes Angebot lediglich einen regulären Arbeitsvertrag erfordert. Vielen Dank für die Prüfung meiner Bewerbung.

\closing{Mit freundlichen Grüßen,}{Radheem Bin Razi}

\end{document}
